% fig62-segtree.tex -- Figure 12.x: segment tree over 8 leaves; node v has children 2v, 2v+1.
% fstyle.tex -- shared TikZ styles for the figures of this book (palette from lifeos.sty).
\tikzset{%
  cell/.style={draw=ink, minimum width=7mm, minimum height=6.5mm, inner sep=0pt, font=\footnotesize, fill=white},
  cellh/.style={cell, fill=accentsoft, draw=accent},
  cella/.style={cell, fill=amberbg, draw=amber},
  cellg/.style={cell, fill=goodbg, draw=good},
  cellx/.style={cell, fill=graybg, draw=slate, text=slate},
  vtx/.style={circle, draw=ink, minimum size=6.5mm, inner sep=0pt, font=\footnotesize, fill=white},
  vtxh/.style={vtx, draw=accent, fill=accentsoft},
  vtxa/.style={vtx, draw=amber, fill=amberbg},
  vtxg/.style={vtx, draw=good, fill=goodbg},
  vtxd/.style={vtx, draw=slate, fill=graybg, text=slate},
  lab/.style={font=\footnotesize, text=slate},
  labA/.style={font=\footnotesize, text=accent},
  labM/.style={font=\footnotesize, text=amber},
  idx/.style={font=\tiny, text=slate},
  arr/.style={-{stealth}, thick, ink},
  arrA/.style={-{stealth}, thick, accent},
  arrM/.style={-{stealth}, thick, amber},
  arrG/.style={-{stealth}, thick, good},
  edge/.style={thick, ink},
  edgeA/.style={thick, accent},
  edgeM/.style={thick, amber},
}%
% \cellrow{prefix}{x0}{y}{style}{v0,v1,...}: a row of cells named prefix0, prefix1, ... (pitch 0.75cm)
\providecommand{\cellrow}[5]{\foreach \v [count=\i from 0] in {#5} {\node[#4] (#1\i) at (#2+\i*0.75,#3) {\v};}}%
% \cellidx{prefix}{n}{start}: index labels under cells prefix0 .. prefix(n-1)
\providecommand{\cellidx}[3]{\foreach \i in {0,...,\numexpr#2-1\relax} {\node[idx, below=0pt of #1\i] {\the\numexpr\i+#3\relax};}}%
% \fitfig{<tikzpicture>}: shrink the picture to \linewidth only when it would overflow.
\ifdefined\fitbox\else\newsavebox{\fitbox}\fi%
\providecommand{\fitfig}[1]{\sbox{\fitbox}{#1}\ifdim\wd\fitbox>\linewidth\resizebox{\linewidth}{!}{\usebox{\fitbox}}\else\usebox{\fitbox}\fi}%

\begin{tikzpicture}[font=\footnotesize]
\tikzset{sn/.style={draw=slate, fill=accentsoft, rounded corners=1pt, minimum height=7mm, minimum width=1.3cm, inner sep=1pt, align=center, font=\scriptsize}}
\foreach \v/\lo/\hi/\x/\y in {1/1/8/5.25/3.0, 2/1/4/2.625/2.0, 3/5/8/7.875/2.0, 4/1/2/1.3125/1.0, 5/3/4/3.9375/1.0, 6/5/6/6.5625/1.0, 7/7/8/9.1875/1.0}
  \node[sn] (n\v) at (\x*1.0,\y*1.0) {{[\lo,\hi]}\\[-1pt]{\tiny node \v}};
\foreach \v/\lo/\x in {1/1/0.65, 2/2/1.97, 3/3/3.28, 4/4/4.6, 5/5/5.9, 6/6/7.2, 7/7/8.5, 8/8/9.8}
  \node[cell, minimum width=1.0cm, minimum height=6mm, font=\tiny] (l\v) at (\x*1.0+0.0,0) {$a_\lo$};
\draw[edge] (n1)--(n2) (n1)--(n3) (n2)--(n4) (n2)--(n5) (n3)--(n6) (n3)--(n7);
\draw[edge] (n4)--(l1) (n4)--(l2) (n5)--(l3) (n5)--(l4) (n6)--(l5) (n6)--(l6) (n7)--(l7) (n7)--(l8);
\node[lab, anchor=west] at (0,-0.9) {array of size $2n$: node $v$ stores the combine of its segment; children are $2v$ and $2v{+}1$};
\end{tikzpicture}
